\section{Modelagem do transporte intramunicipal de vacinas}

A distribuição intramunicipal de vacinas corresponde à etapa final da cadeia logística, na qual as doses são transportadas da secretaria municipal de saúde para os pontos de aplicação, como postos de saúde e hospitais, ou seja, existe apenas um ponto de origem das vacinas.

Para o caso intramunicipal, somente o modal rodoviário é considerado.

\subsection{Variável de Decisão e Parâmetros}

Seja:

\[
z_{jpv} \in \mathbb{Z}_{\geq 0}
\]

onde:
\begin{itemize}
    \item $z_{jpv}$ representa a quantidade de vacinas do tipo $v$ transportadas do centro do município $j$ para o posto de destino $b$;
    \item $j \in J$ representa os municípios;
    \item $p \in P_j$ representa os postos de saúde pertencentes ao município $j$;
    \item $v \in V$ representa os tipos de vacina.
\end{itemize}

Como se trata de unidades discretas (doses), temos:

\begin{equation*}
z_{jpv} \in \mathbb{Z}_{\geq 0}, \quad \forall j \in J, p \in P_j, v \in V
\end{equation*}

De forma análoga ao modelo intermunicipal, temos:

\[
\tau_{jp}^{rod} \in \mathbb{R}_{\geq 0}
\]
\[
\delta^{rod} \in \mathbb{R}_{\geq 0}
\]

onde:
\begin{itemize}
    \item $\tau_{jp}^{rod}$ é o tempo de deslocamento do modal rodoviário entre o centro de distribuição municipal $j$ e o posto de destino $b$.
    \item $\delta^{rod}$ representa o tempo de processamento logístico do modal rodoviário.
\end{itemize}

Seja:

\begin{equation*}
C_{jp} \in \mathbb{R}_{\geq 0}
\end{equation*}

o \textbf{Custo Total do transporte} entre os pontos $j$ e $b$.

Dessa forma, o \textbf{Custo Total de transporte} é dado por:

\[
C_{jp} = \tau_{jp}^{rod} + \delta^{rod}
\]

No contexto intramunicipal, considera-se que o transporte é realizado exclusivamente por meio do modal rodoviário. Dessa forma, a capacidade de transporte não depende da escolha entre diferentes modais.

Seja:

\[
q_{jp} \in \mathbb{R}_{\geq 0}
\]

onde:
\begin{itemize}
    \item $q_{jp}$ representa a capacidade máxima de transporte entre o centro de distribuição do município $j$ e o posto de destino $p$.
\end{itemize}

e o \textbf{tempo de distribuição} das vacinas é:

\[
T = \max \left( \quad C_{jp} \quad \mid \quad z_{jpv} > 0 \quad \right)
\]

\subsection{Função Objetivo}

A \textbf{função objetivo} é definida como:

\begin{equation}
\min z = \sum_{v} \sum_{j} \sum_{p \in P_j} C_{jp} \cdot z_{jpv}
\label{intramunicipal_sem_penalidade}
\end{equation}

O objetivo é minimizar o tempo total de distribuição das vacinas dentro do município.

Para incorporar o risco de expiração das vacinas, introduz-se um fator de penalização associado ao tempo restante de validade no município. Seja:

\[
\rho_{jv} = \frac{1}{e_{jv}}
\]
onde:
\begin{itemize}
\item  $e_{jv}$ representa o tempo restante até a expiração da vacina do tipo $v$ disponível no município $j$.
\item $\rho_{jv}$ é penalidade da vacina do tipo $v$ no município $j$
\end{itemize}

Dessa forma, a função objetivo passa a ser:
\begin{equation}
\min z = \sum_{v} \sum_{j} \sum_{p \in P_j} C_{jp} \cdot z_{jpv} \cdot \rho_{jv}
\label{intramunicipal_com_penalidade}
\end{equation}

Essa formulação prioriza a distribuição de vacinas com menor tempo restante de validade associado ao estoque no nível intramunicipal.

\subsection{Restrições}

A \textbf{restrição de oferta} diz respeito a quantidade total enviada por cada ponto de origem não poder exceder sua disponibilidade:

\begin{equation*}
\sum_{p \in P_j} z_{jpv} \le s_{jv}, \quad \forall j \in J, \; v \in V
\end{equation*}

onde:
\begin{itemize}
    \item $s_{jv}$ representa a quantidade disponível da vacina $t$ no ponto de origem $j$.
\end{itemize}

A \textbf{restrição de demanda} estabelece que cada ponto de destino deve receber exatamente sua demanda:

\begin{equation*}
\sum_{j} z_{jpv} = r_{pv}, \quad \forall p \in P_j, \; v \in V
\end{equation*}

onde:
\begin{itemize}
    \item $r_{pv}$ representa a demanda da vacina $t$ no ponto de destino $p$.
\end{itemize}

A \textbf{restrição de capacidade} considera que a capacidade de transporte entre os pontos é limitada:

\begin{equation*}
\sum_{v} z_{jpv} \le q_{jp}, \quad \forall j \in J, p \in P_j
\end{equation*}

onde:
\begin{itemize}
    \item $q_{jp}$ representa a capacidade máxima de transporte entre $j$ e $p$.
\end{itemize}

A \textbf{restrição de viabilidade} temporal é definida para evitar perdas por expiração:

\begin{equation*}
z_{jpv} = 0, \quad \forall j \in J, \; p \in P_j, \; v \in V \text{ tais que } C_{jp} > e_{jv}
\end{equation*}

onde:
\begin{itemize}
    \item $e_{jv}$ representa o tempo restante até a expiração da vacina $v$ no ponto de origem $a$.
\end{itemize}

O modelo completo é dado por:

\begin{equation*}
\min z = \sum_{v} \sum_{j} \sum_{b} C_{jp} \cdot z_{jpv}
\end{equation*}
sujeito a:
\begin{equation*}
\sum_{b} z_{jpv} \le s_{jv}, \quad \forall j \in J, \; v \in V
\end{equation*}
\begin{equation*}
\sum_{j} z_{jpv} = r_{pv}, \quad \forall p \in P_j, \; v \in V
\end{equation*}
\begin{equation*}
\sum_{v} z_{jpv} \le q_{jp}, \quad \forall j \in J, p \in P_j
\end{equation*}
\begin{equation*}
z_{jpv} = 0, \quad \forall j \in J, \; p \in P_j, \; v \in V \text{ tais que } C_{jp} > e_{jv}
\end{equation*}
\begin{equation*}
z_{jpv} \in \mathbb{Z}_{\geq 0}, \quad \forall j \in J, p \in P_j, v \in V
\end{equation*}


\bigskip